\subsection{逼近性质}

引理 1. --- 设 X 是一个局部紧空间，K 是 X 的一个紧子集，\((\mathrm{V}_{k})_{1\leqslant k\leqslant n}\) 是 K 由 X 的开集组成的一个有限覆盖。则存在 n 个连续映射 \(f_{k}\) （把 X 映到 {[}0,1{]} 中），使得 \(f_{k}\) 的支集含于 \(V_{k}\) （对 \(1\leqslant k\leqslant n\)），且 \(\sum_{k=1}^{n}f_{k}(x)\leqslant1\) 对所有 \(x\in X\) 成立，\(\sum_{k=1}^{n}f_{k}(x)=1\) 对所有 \(x\in K\) 成立。

因为，设 \(X'\) 是由 X 添加一个无穷远点 \(\omega\) 而得到的紧空间（GT, I, §9, No. 8, Th. 4）；诸集合 \(V_{0} = X' - K\) 与 \(V_{k}\) （ \(1 \leqslant k \leqslant n\) ）构成 \(X'\) 的一个开覆盖。设 \((f_{k})_{0 \leqslant k \leqslant n}\) 是从属于 \(X'\) 的这一覆盖的一个连续单位分解（GT, IX, §4, No. 3, 命题 3）；诸函数 \(f_{k}\) （指标 \(k \geqslant 1\) ）满足引理的条件。

引理 2. --- 设 X 是一个局部紧空间，K 是 X 的一个紧子集，E 是一个局部凸空间，q 是 E 上的一个连续半范数，\(\Phi\) 是 X 到 E 中的映射所成的一个等度连续集，其诸支集都含于 K。则对每个 \(\varepsilon > 0\) ，存在 X 上的一个连续单位分解 \((\varphi_{j})_{0 \leqslant j \leqslant n}\) ，具有下列性质：

\begin{enumerate}
\def\labelenumi{(\roman{enumi})}
\item
  \(\operatorname{Supp}(\varphi_j) \subset \mathrm{K}\) 对 \(1 \leqslant j \leqslant n\) 成立。
\item
  若对 \(1 \leqslant j \leqslant n\) ，\(x_j\) 是 \(\operatorname{Supp}(\varphi_j)\) 的任意一点，则对每个函数 \(\mathbf{f} \in \Phi\) 及每个 \(x \in \mathbf{X}\) ，
\end{enumerate}

\[
q \left(\mathbf {f} (x) - \sum_ {j = 1} ^ {n} \varphi_ {j} (x) \mathbf {f} \left(x _ {j}\right)\right) \leqslant \varepsilon .\tag{1}
\]

对属于 K 的边界的每个 \(y\) ，有 \(\mathbf{f}(y) = 0\) （对一切 \(\mathbf{f} \in \Phi\) ），因此存在一个开邻域 \(\mathrm{V}_y\) （\(y\) 在 X 中的），使得对每个 \(z \in \mathrm{V}_y\) 及每个 \(\mathbf{f} \in \Phi\) ，有 \(q(\mathbf{f}(z)) \leqslant \varepsilon / 2\) 。设 \(\mathrm{K}'\) 是 K 中不属于任何 \(\mathrm{V}_y\) （当 \(y\) 取遍 K 的边界）的点所成的集；\(\mathrm{K}'\) 是紧的且含于 K 的内部。集 \(\Phi\) 在 K 中是一致等度连续的；因此存在一个有限开覆盖 \((\mathrm{U}_j)_{1 \leqslant j \leqslant n}\) （\(\mathrm{K}'\) 的，由含于 K 的 X 的开集组成），使得对每对点 \(x, y\) （同一个 \(\mathrm{U}_j\) 中的），有 \(q(\mathbf{f}(x) - \mathbf{f}(y)) \leqslant \varepsilon / 2\) （对每个 \(\mathbf{f} \in \Phi\) ）。由引理 1，存在 \(n\) 个连续映射 \(\varphi_j\) （X 到 {[}0,1{]} 中的； \(1 \leqslant j \leqslant n\) ），使得 \(\operatorname{Supp}(\varphi_j) \subset \mathrm{U}_j\) ，且 \(\sum_{j=1}^{n} \varphi_j(x) \leqslant 1\) 在 X 上成立，\(\sum_{j=1}^{n} \varphi_j(x) = 1\) 在 \(\mathrm{K}'\) 上成立。对 \(x_j \in \operatorname{Supp}(\varphi_j)\) （ \(1 \leqslant j \leqslant n\) ）及 \(\mathbf{f} \in \Phi\) ，我们于是有：对一切 \(x \in \mathrm{U}_j\) ，

\[
q \big (\mathbf {f} (x) \varphi_ {j} (x) - \mathbf {f} (x _ {j}) \varphi_ {j} (x) \big) = \varphi_ {j} (x) q \big (\mathbf {f} (x) - \mathbf {f} (x _ {j}) \big) \leqslant \frac {\varepsilon}{2} \varphi_ {j} (x),
\]

且当 \(x \notin U_{j}\) 时这一关系仍成立，因为这时 \(\varphi_{j}(x) = 0\) 。相加即得：对每个 \(x \in X\) ，

\[
q \left(\mathbf {f} (x) \left(1 - \varphi_ {0} (x)\right) - \sum_ {j = 1} ^ {n} \varphi_ {j} (x) \mathbf {f} \left(x _ {j}\right)\right) \leqslant \frac {\varepsilon}{2} \left(1 - \varphi_ {0} (x)\right),\tag{2}
\]

其中 \(\varphi_0 = 1 - \sum_{j=1}^{n} \varphi_j\) ；由此当 \(x \in \mathbf{K}'\) 时得 (1)，因为这时 \(\varphi_0(x) = 0\) ；当 \(x \notin \mathbf{K}\) 时 (1) 也成立，因为这时第一项为零。最后，对 \(x \in \mathbf{K} - \mathbf{K}'\) ，有 \(q(\mathbf{f}(x)\varphi_0(x)) \leqslant \varepsilon/2\) （由 \(\mathbf{K}'\) 的定义），因此这一关系与 (2) 在此情形仍蕴涵 (1)。

设 X 是一个局部紧空间；对每个 Banach 空间 E（实的或复的），我们用 \(\mathcal{C}^{b}(X;E)\) 表示 X 到 E 中的连续且有界的映射所成的向量空间；已知 X 中的一致收敛拓扑与 \(\mathcal{C}^{b}(X;E)\) 的（实的，相应地，复的）向量空间结构相容，且它由范数

\[
\| \mathbf {f} \| = \sup _ {x \in X} \| \mathbf {f} (x) \|.\tag{3}
\]

此外，这样定义的赋范空间是一个 Banach 空间（GT, X, §3, No. 2, 以及 No. 1, 命题 2 的推论 2）；这一范数在 \(\mathcal{K}(\mathrm{X};\mathrm{E})\) 上定义的拓扑（换句话说，X 中的一致收敛拓扑）粗于第 1 目中在 \(\mathcal{K}(\mathrm{X};\mathrm{E})\) 上定义的直接极限拓扑。

命题 3. --- 设 X 是一个局部紧空间，X‘ 是由 X 添加一个无穷远点 ω 而得到的紧空间（GT, I, §9, No. 8,

Th. 4)，E 是一个 Banach 空间。\(\mathcal{K}(\mathrm{X};\mathrm{E})\) 在赋范空间 \(\mathcal{C}^b (\mathrm{X};\mathrm{E})\) 中的闭包是 \(\mathbf{X}\) 上以 E 为值且在点 \(\omega\) 处趋于 0 的连续函数所成的向量空间。

设 \(\mathbf{f} \in \mathcal{C}^b(\mathrm{X};\mathrm{E})\) 是属于 \(\mathcal{K}(\mathrm{X};\mathrm{E})\) 的闭包中的一个函数；对每个 \(\varepsilon > 0\) ，存在一个函数 \(\mathbf{g} \in \mathcal{K}(\mathrm{X};\mathrm{E})\) ，使得 \(\| \mathbf{f}(x) - \mathbf{g}(x) \| \leqslant \varepsilon\) 对所有 \(x \in \mathrm{X}\) 成立；若 \(\mathrm{K}\) 是 \(\mathbf{g}\) 的支集，则由此得出 \(\| \mathbf{f}(x) \| \leqslant \varepsilon\) 对所有 \(x \in \complement \mathrm{K}\) 成立，从而 \(\mathbf{f}(x)\) 当 \(x\) 趋于 \(\omega\) 时趋于 0。反之，若 \(\mathbf{f}\) 具有这一性质，则对每个 \(\varepsilon > 0\) ，存在一个紧集 \(\mathrm{K} \subset \mathrm{X}\) ，使得 \(\| \mathbf{f}(x) \| \leqslant \varepsilon\) 对所有 \(x \in \complement \mathrm{K}\) 成立。由引理 1，存在一个连续映射 \(h\) （X 到 {[}0,1{]} 中的），具有紧支集且在 \(\mathrm{K}\) 上等于 1；这时 \(\| \mathbf{f}(x) h(x) \| \leqslant \varepsilon\) 在 \(\complement \mathrm{K}\) 上成立，且 \(\mathbf{f}(x) = \mathbf{f}(x) h(x)\) 在 \(\mathrm{K}\) 上成立；由于 \(\mathbf{f}h\) 具有紧支集且 \(\| \mathbf{f}(x) - \mathbf{f}(x) h(x) \| \leqslant 2\varepsilon\) 对所有 \(x \in \mathrm{X}\) 成立，命题得证。

我们将用 \(\mathcal{C}^{0}(\mathrm{X};\mathrm{E})\) 表示 \(\mathcal{C}^{b}(\mathrm{X};\mathrm{E})\) 中由在无穷远点 \(\omega\) 处趋于零的函数组成的子空间；它也就是赋范空间 \(\mathcal{K}(\mathrm{X};\mathrm{E})\) 的完备化。

命题 4. --- 设 X 是一个局部紧空间，E 是一个局部凸空间；则空间 \(\mathcal{K}(\mathrm{X};\mathrm{E})\) 在 \(\mathcal{C}(\mathrm{X};\mathrm{E})\) 中对紧收敛拓扑是稠密的。

对每个紧集 \(K \subset X\) ，由引理 1，存在一个在 K 上等于 1 的函数 \(h \in \mathcal{K}(X; \mathbf{R})\) ；对每个函数 \(\mathbf{f} \in \mathcal{C}(X; E)\) ，函数 hf 属于 \(\mathcal{K}(X; E)\) 且在 K 上等于 f，由此得我们的断言。

命题 5. --- 设 X 是一个局部紧空间，E 是一个实的（相应地，复的）局部凸空间。对 X 的每个紧子集 K，向量空间 \(\mathcal{K}(X,K;\mathbf{R})\otimes_{\mathbf{R}}E\) （相应地，\(\mathcal{K}(X,K;\mathbf{C})\otimes_{\mathbf{C}}E\) ）（等同于 X 到 E 中的映射所成的一个集，参见 A, II, §7, No. 7, 命题 15 的推论）在 \(\mathcal{K}(X,K;E)\) 中稠密；向量空间 \(\mathcal{K}(X;\mathbf{R})\otimes_{\mathbf{R}}E\) （相应地，\(\mathcal{K}(X;\mathbf{C})\otimes_{\mathbf{C}}E\) ）在 \(\mathcal{K}(X;E)\) 中稠密。

由于第二个断言是第一个断言的明显推论，只需证明后者。我们应用引理 2，其中 \(\Phi\) 化为 \(\mathcal{K}(X,K;E)\) 的单个元素 f；这时，对每个 \(x\in X\) ，

\[
q \left(f (x) - \sum_ {j = 1} ^ {n} \varphi_ {j} (x) f \left(x _ {j}\right)\right) \leqslant \varepsilon ,
\]

其中 \(\varphi_{j}\) 属于 \(\mathcal{K}(\mathrm{X},\mathrm{K};\mathbf{R})\) ；由于映射 \(x\mapsto \sum_{j = 1}^{n}\varphi_{j}(x)f(x_{j})\) 可以典范地等同于元素 \(\sum_{j = 1}^{n}\varphi_{j}\otimes f(x_{j})\) ，由 \(\mathcal{K}(\mathrm{X},\mathrm{K};\mathrm{E})\) 的拓扑的定义，这就证明了命题。

